\section{On discovering you're not an Android}

\begin{quotex}
The Turing test was developed to expose androids masquerading as humans. But what if humans themselves are actually androids? 
\end{quotex}

Over at Gene Expressions\footnote{\url{https://www.gnxp.com/WordPress/2011/07/04/on-discovering-youre-an-android/}}, we learn this alarming fact:

\begin{quotex}
The idea that the self, or the conscious mind, emerges from the workings of the physical structures of the brain — with no need to invoke any supernatural spirit, essence or soul — is so fundamental to modern neuroscience that it almost goes unmentioned. It is the tacitly assumed starting point for discussions between neuroscientists, justified by the fact that all the data in neuroscience are consistent with it being true. Yet it is not an idea that the vast majority of the population is at all comfortable with or remotely convinced by. Its implications are profound and deeply unsettling, prompting us to question every aspect of our most deeply held beliefs and intuitions.

\end{quotex}
Although I have great respect for the intelligence and learning of scientists, they persist in being philosophically naive. Materialism leads to empiricism which leads to Kantian scepticism which leads to idealism, the very opposite of the starting point. A methodological assumption that is appropriate for neuroscientists is not tantamount to a proof of its universal validity.

Astronomy, for example, assumed that the motion of the planets could be explained by matter and the force of gravity. That assumption worked very well and astronomers built up a large inventory of observations of planetary motion. Unexpected anomalies in those observations led to the discovery of more planets beyond Saturn, which would have been inexplicable had astronomers assumed there were only six planets. Neuroscientists are at the incipient stage of building a library of observations. There is a lot to do if they hope to explain all subjective inner states of consciousness by means of material electrochemical processes. But that is the first problem they face: they do not accept the reality of inner states of consciousness, but rely on what the experimental subject reports.

\begin{quotex}
all those kinds of experiments have, of course, been done on human beings, tens of thousands of times. Functional magnetic resonance imaging methods let us correlate the activity of particular brain circuits with particular behaviours, perceptions or reports of inward states.

\end{quotex}
Since inner states cannot be directly observed by any apparatus, the scientist relies on the ``report" of the experimental subject. Let's say, for example, that when his neurons are stimulated a certain way, the subject reports, ``I felt a pleasant sensation accompanied by a vision of yellowish colours." Does the scientist then take out his oscilloscope to analyze the report in a scientific manner? Of course not. He relies on his own subjective experience of pleasantness and yellow to imagine what the subject reported. On the assumptions of the neuroscientist, we must conclude that when the subject's neurons are stimulated in a certain way, the subject will then emit sound waves of a certain frequencies and duration, which will then stimulate the corresponding neurons in the brain of the neuroscientist. Maybe for Occam this seems like the simplest explanation, but for the rest of us, something rather essential seems to be left out.

\textbf{A second difficulty} is that, for the neuroscientist, inner states don't really exist. This is truly awkward, since his goal is to explain the experience of inner states by bio-electro-chemical processes. But we are invited to participate in this thought experiment:

\begin{quotex}
Imagine you came across a robot that performed all the functions a human can perform – that reported a subjective experience apparently as rich as yours. If you were able to observe that the activity of certain circuits was associated with the robot's report of subjective experience, if you could drive that experience by activating particular circuits, if you could alter it by modifying the structure or function of different circuits, would there be any doubt that the experience arose from the activity of the circuits?

\end{quotex}
This presumably demonstrates that inner experiences are not necessary to produce a ``report" of them. However, the human subject believes he is reporting an actual experience. The usual solution is to call it an \emph{epiphenomenon}, that is, a meaningless side effect of the bio-electro-chemical process, much like our experience of the rising and setting of the sun. The epiphenomenon cannot exist, otherwise it would have effects. But the scientist rules that out from the start.

\textbf{A third problem} is the notion of truth. Material processes are neither true nor false. Those are logical, not physical, categories of thought and apply only to propositions. So our neuroscientists observes the effects of his experiments and concludes with the theory: ``All inner states of the mind are the result of material processes in the brain." He wants to convince us all that this is true.

Presumably, since a thought is an inner state, the report of a thought is the result of a material process. Let's call the neuroscientist's theory proposition A. So when I have a thought of A, that is the result of a material process. Yet so is the thought of not-A. So how do I determine which is true? That requires another thought: ``A is true", which is the result of a different process, as is its opposite, ``not-A is true". Hence, the scientist can stimulate my neurons and make me believe either proposition. There is no necessity to convince me by arguments. Obviously, an argument is a report that stimulates the neurons of the hearer to believe the argument, at least in those brains trained to ``understand" the argument, whatever that may mean. But very few have the intelligence and resources to become neuroscientists. Hence, we can shortcut the process and just provide them with a neuron stimulation kit. We could then provide neuron simulation kits to make users experience or believe all sorts of things. We can make everyone equally intelligent, since they would then know all true things (given the right kits), or courageous or just and so on.

\begin{quotex}
I will leave the neuroscientist with another thought experiment. He will hook the human experimental subject up to his apparatus. Before each experiment, he will tell the subject what the subject will report. The scientist will then stimulate the subject's brain and wait for his report. Eventually, the subject will be convinced that the scientist's theory is correct since he is able to anticipate the subject's reports in advance. He will, won't he?
\end{quotex}


\flrightit{Posted on 2011-12-28 by Cologero }

\begin{center}* * *\end{center}

\begin{footnotesize}\begin{sffamily}



\texttt{Golgonooza on 2011-12-28 at 05:22 said: }

Again, excellent post Cologero. Another critique of the reductive materialist or functionalist theory of mind comes from Searle's Chinese Room argument. Imagine someone who knows only English sitting in a room alone with a rule book in English for manipulating strings of Chinese characters. Someone who passes a question into the room from outside in Chinese will get an appropriate answer back in Chinese and it will look to them as if the person in the room understands it. But inside the room, the person is just using a rule book to construct the characters and understands no Chinese. It is often used as an argument against the possibility of AI. But also it implies that meaning is not equivalent to the sum total of physical processes that convey it.


\hfill

\texttt{Cologero on 2011-12-29 at 09:35 said: }

It makes you wonder, Golgonooza, about the inner states of awareness of some people. Perhaps their consciousnesses are rather bland so their theories seem plausible. Tomberg gives the example of the materialist reading a book. He believes the letters wrote themselves, then combined themselves into syllables as the result of chemical reactions of the ink. The letters and words are then nothing but epiphenomena of the ink.

Yet this way of thinking is considered cutting edge and its adherents call themselves ``bright".

Ultimately it is the result of an incorrect understanding of the Traditional notions of syntax and semantics. The materialist only accepts syntax and denies semantics. Unfortunately, that denial is not possible in practice, so semantics always creeps in … words actually ``mean" something.

This is a throwback to Hilbert's challenge to produce an automaton that proves mathematical theorems. In effect, mathematics must be reduced to syntax alone, without semantic meaning. Godel showed that no such automaton is possible, but news travels slowly. Beyond that, of course, materialism is a continuation of Bacon's project of eliminating formal and final causes. In this view, phenomena are what they are and have no meaning or objectivity beyond that. It is successful as a methodological assumption but not as a guide to life.


\hfill

\texttt{Attila on 2011-12-29 at 10:53 said: }

Any system or mechanism will always be less complex than its creator.

Perhaps there is a connection between homosexuality and a this-is-all-there-is view of life – which could shed some light on Bacon's alleged condition and his ``philosophy" – abusing the term here. 


\hfill

\texttt{Cologero on 2012-02-20 at 20:42 said: }

Beyond the philosophical issue, I proposed what I believe to be a scientific test. I don't know of any such test proposed by a neuroscientist; if there is one, I would like to know.


\hfill

\texttt{Sparrow on 2015-10-23 at 08:58 said: }

When one considers the opposite of the materialist proposition: that psyche exists before matter and puts it into motion, not the other way around, everything falls into place. The materialist should seriously consider this, as it better explains their worldview better than their current proposition that single celled organisms inevitably evolved to have brain cells and produce consciousness. Of course, their beliefs would still be flawed, but at least a little less so.


\hfill

\texttt{White Rabbit on 2015-10-23 at 14:41 said: }

I've never been a materialist but the idea that psyche is primary has always seemed counter-intuitive to me. Whether you take psychoactives or bump your head, your psyche follows matter. That needs to be accounted for.


\hfill

\texttt{obscure on 2015-10-24 at 01:24 said: }

White Rabbit,

What most people mean by consciousness is usually too broad. Obviously one is conscious of one's body and is thus a subject receiving various internal and external bodily modifications. It would therefore be best to consider this broad definition of consciousness and then strip away everything. The only thing left is the purest intentionality; the external efficiency of which we may call will. But, what is of primary interest is simply the pure, intellectual intention which must even exclude the simplest external efficiency. Remind yourself of Plotinus' spiritual exercises involving the imaginary sphere. These types of exercises are so clear, simple and necessary for removing all of the unnecessary associations usually attached to the pure intellect. In the purest intention the concept intended and the agent intending are identical; the concept intended may be called `Absolute Unity' or `Universal Possibility'. The supposed problems presented by bodily modifications ought to disappear here and also the superficiality of most psychedelic drug use ought to be very apparent.

Nonetheless, it will also be clear that everything which is experienced is also a modification of Universal Possibility and therefore it is possible to subsume even the grossest phenomena we encounter through a correct understanding. But, he who attains the highest states experiences everything in the highest possible way and thus there is no need to depart from mundane phenomena in one's external activity. Indeed, everything experienced externally (including the internal sensation of the body and its affections) insofar as it is not also understood to signify something beyond itself (and is thus only taken as an external ejaculation of physical force; subjected to decay and disappearance) is absolutely mundane. So, the spiritual man does not see the mundane things because he only sees signs and so his criterion does not direct him to strange phenomena, but only to his tasks.


\hfill

\texttt{White Rabbit on 2015-10-24 at 14:25 said: }

I concur Obscure, but when we strip away everything we strip away the particular ``inner states" Cologero referred to. This is what I meant by psyche.


\hfill

\texttt{obscure on 2015-10-25 at 02:32 said: }

The purpose of stripping away everything is to see the real ground for all of those later consequences so that one does not have crude pictorial thoughts about matter. The materialist thinks that he has the same plan in his research, but he is moving towards the periphery rather than the center of things. 

To re-state this topic in another way: One must strip away the inner states in order to grasp the spiritual ground in the same way that the materialist believes that one must strip away those states in order to see small bodies impinging upon one another (or more probably the physical models of such things). The point is that he does not reach the term, but he proceeds indefinitely to collect data and conventionally organize it. His activity is justified by technological progress. Our activity, on the other hand, is only justified by spiritual experience which is necessarily less egalitarian, less blatant, more difficult, more obscure, etc. than what is readily perceived in the techno-culture. So, only certain presentations of metaphysical arguments or valid criticisms of the immaturity of materialism can suffice from an exoteric standpoint. Under this latter aspect we see the character of Cologero's post here which has nothing to do with the spiritual `stripping away' of inner states. Sometimes one may indicate the right direction not so much by a definite affirmation but by a negation of the contrary.


\hfill

\texttt{Alistair Fraser on 2015-10-25 at 21:37 said: }

The chinese room argument is rationally revealing of the contradictions

Also The proposed scientific test would never be entertained = Why would the WHO expose its own limitations of understanding which is one of the most lucrative evolving ``industry" in modern times.

On neuron stimulations accuracy of anything, a simple mechanical car has electrical, exhaust, fuel circuits , not to mention the hardware and software , of which a fault can develop in any that effects the whole , what makes them think humans are less complex than a simple car ?

Heres something else to think about = Today i watched an excerpt of an experiment in japan whereby a human wore one of those virtual viewers in which they could see through it to the outside world plus one more ingredient was that the computer it was attached too could input virtual influence.

So the computer simulated to the person wearing the viewer that a small green ball was sitting in the air in front of the viewer , the person then lifted their hand up to feel the ball and they confirmed they squeezed and felt it .

After, it was disclosed that the computer had actually re-represented their hand as it moved up to look like it was squeezing the green ball, but what had happened in reality was that the person had simply squeezed / rubbed their fingers together being fooled into believing it was a virtual green ball .

Now what this proved was that , as some recent studies had shown, people give priority of belief to what they see , over that of which they feel .

And i think this is a result of the electronic technological mediated world , where people have been mass-conditioned by the power of television and film to form their opinions so that their instinctive feeling power has been quietly and efficiently subordinated below their seeing faculty , just from sheer laziness , if its shown enough then it must be true , or if its shown on an official platform then it must be true

This then breeds a level of trance-induced imported certainty of opinion that cannot justify itself, and if questioned, reacts with disturbed anger at the questioner as if they be mad


\hfill

\texttt{Cologero on 2015-10-26 at 20:35 said: }

@White Rabbit, who wrote: ``Whether you take psychoactives or bump your head, your psyche follows matter."

That's great news WR. Perhaps you can take a drug to bring you peace, joy, happiness, and enlightenment.


\hfill

\texttt{White Rabbit on 2015-10-26 at 21:20 said: }

That would be possible if supra-human states followed matter like any ordinary states of consciousness, Cologero. Or instead, what Guénon says about symbols, rites, ect. could apply:

``But some will ask, how is it possible that merely contingent means can produce an effect which immeasurably surpasses them and which is of a totally different order from that to which the instruments themselves belong? We should first point out that these means are, in reality, only fortuitous. The results they help to attain are by no means consequential; they place the being in the position requisite for attainment and that is all…

As for us, we do not confuse a simple means with a cause in the true sense of the word and we do not regard metaphysical realization as an effect, since it is not the production of something which does not yet exist, but the knowing of that which is, in an abiding and immutable manner, beyond all temporal succession, for all states of the being, considered under their primary aspect, abide in perfect simultaneousness in the eternal now."


\hfill


\end{sffamily}\end{footnotesize}
